% GENERATED FILE. Edit canonical lessons, then run npm run generate:books.
% canonical-source-hash: fnv1a64:982bf683ad5948e3
% canonical-lessons: RU-C503-kollektsionirovat, RU-C503-kollektsioner, RU-C503-znachok, RU-C503-krossvord, RU-C503-razgadyvat

\chapter{Collections and Puzzles}
\label{ch:ru-collections-and-puzzles}
\hlchaptermodality{503}

\begin{chapteropening}
\textbf{By the end of this chapter:} \emph{I can talk about collecting things and doing puzzles.}

Complete the last lesson of chapter 503: I can talk about collecting things and doing puzzles.
\end{chapteropening}

\section[\texorpdfstring{\ru{коллекционировать} (kollektsionírovat)}{коллекционировать (kollektsionírovat)}]{\ru{коллекционировать} (kollektsionírovat) — to collect (as a hobby)}

\label{lesson:RU-C503-kollektsionirovat}



\begin{quote}
Before the new one: say the Russian for to make things by hand, then the Russian for a handmade craft item.
\end{quote}

\subsection*{You'll want to know: \ru{коллекционировать}}
\textbf{\ru{коллекционировать}} (\emph{kollektsionírovat}) — \textquotedblleft{}to collect (as a hobby)\textquotedblright{}.

\begin{quote}
\ru{Мой} \ru{сосед} \ru{коллекционирует} \ru{старые} \ru{монеты}.
\end{quote}

\begin{quote}
\emph{Moy soséd kollektsioníruyet stáryye monéty.}
\end{quote}

\begin{quote}
My neighbour collects old coins.
\end{quote}

Imperfective; present \ru{коллекционирую}.

\begin{cousinweb}
\emph{\ru{коллекционировать}} is built on \emph{\ru{коллекция}} 'a collection', a learned borrowing usually traced to Latin \emph{collēctiō} 'a collecting together' (Lewis \& Short), from \emph{colligere}; EtymDB lists the verb without a source.
\end{cousinweb}

\begin{grammarlens}[title={What you've built}]
One more word for telling, explaining and describing at B1.
\end{grammarlens}

\subsection*{Your turn}
\begin{itemize}
\raggedright
  \item \emph{Say it:} \emph{kollektsionírovat}
  \item \emph{Say it:} \emph{kollektsionírovat}, once more
  \item \emph{Say it:} a sentence of your own with \emph{kollektsionírovat}
  \item \emph{Recall:} say the Russian for to make things by hand, then the Russian for a handmade craft item, then say \emph{kollektsionírovat} again
\end{itemize}

\begin{tcolorbox}[breakable,colback=teal!4,colframe=teal!35!black,title={Before you move on}]
What does \ru{коллекционировать} mean? (\textquotedblleft{}To collect (as a hobby)\textquotedblright{}.) Say it once more.
\end{tcolorbox}
\section[\texorpdfstring{\ru{коллекционер} (kollektsionér)}{коллекционер (kollektsionér)}]{\ru{коллекционер} (kollektsionér) — a collector}

\label{lesson:RU-C503-kollektsioner}



\begin{quote}
Before the new one: say the Russian for a handmade craft item, then the Russian for to collect.
\end{quote}

\subsection*{You'll want to know: \ru{коллекционер}}
\textbf{\ru{коллекционер}} (\emph{kollektsionér}) — \textquotedblleft{}a collector\textquotedblright{}.

\begin{quote}
\ru{Известный} \ru{коллекционер} \ru{купил} \ru{эту} \ru{картину}.
\end{quote}

\begin{quote}
\emph{Izvéstnyy kollektsionér kupíl étu kartínu.}
\end{quote}

\begin{quote}
A well-known collector bought this painting.
\end{quote}

Masculine.

\begin{cousinweb}
\emph{\ru{коллекционер}} is built on \emph{\ru{коллекция}} 'collection', usually traced to Latin \emph{collēctiō} 'a collecting together' (Lewis \& Short); EtymDB has no entry.
\end{cousinweb}

\begin{grammarlens}[title={What you've built}]
One more word for telling, explaining and describing at B1.
\end{grammarlens}

\subsection*{Your turn}
\begin{itemize}
\raggedright
  \item \emph{Say it:} \emph{kollektsionér}
  \item \emph{Say it:} \emph{kollektsionér}, once more
  \item \emph{Say it:} a sentence of your own with \emph{kollektsionér}
  \item \emph{Recall:} say the Russian for a handmade craft item, then the Russian for to collect, then say \emph{kollektsionér} again
\end{itemize}

\begin{tcolorbox}[breakable,colback=teal!4,colframe=teal!35!black,title={Before you move on}]
What does \ru{коллекционер} mean? (\textquotedblleft{}A collector\textquotedblright{}.) Say it once more.
\end{tcolorbox}
\section[\texorpdfstring{\ru{значок} (znachók)}{значок (znachók)}]{\ru{значок} (znachók) — a badge}

\label{lesson:RU-C503-znachok}



\begin{quote}
Before the new one: say the Russian for to collect, then the Russian for a collector.
\end{quote}

\subsection*{You'll want to know: \ru{значок}}
\textbf{\ru{значок}} (\emph{znachók}) — \textquotedblleft{}a badge\textquotedblright{}.

\begin{quote}
\ru{У} \ru{него} \ru{на} \ru{рюкзаке} \ru{много} \ru{значков}.
\end{quote}

\begin{quote}
\emph{U negó na ryukzaké mnógo znachkóv.}
\end{quote}

\begin{quote}
He has a lot of badges on his backpack.
\end{quote}

Masculine; genitive \ru{значка}. A diminutive of \emph{\ru{знак}} — a sign.

\begin{cousinweb}
\emph{\ru{значок}} is a diminutive of \emph{\ru{знак}} 'a sign', built on the root of \emph{\ru{знать}} 'to know'; EtymDB has no entry for \emph{\ru{значок}}.
\end{cousinweb}

\begin{grammarlens}[title={What you've built}]
One more word for telling, explaining and describing at B1.
\end{grammarlens}

\subsection*{Your turn}
\begin{itemize}
\raggedright
  \item \emph{Say it:} \emph{znachók}
  \item \emph{Say it:} \emph{znachók}, once more
  \item \emph{Say it:} a sentence of your own with \emph{znachók}
  \item \emph{Recall:} say the Russian for to collect, then the Russian for a collector, then say \emph{znachók} again
\end{itemize}

\begin{tcolorbox}[breakable,colback=teal!4,colframe=teal!35!black,title={Before you move on}]
What does \ru{значок} mean? (\textquotedblleft{}A badge\textquotedblright{}.) Say it once more.
\end{tcolorbox}
\section[\texorpdfstring{\ru{кроссворд} (krossvórd)}{кроссворд (krossvórd)}]{\ru{кроссворд} (krossvórd) — a crossword}

\label{lesson:RU-C503-krossvord}



\begin{quote}
Before the new one: say the Russian for a collector, then the Russian for a badge.
\end{quote}

\subsection*{You'll want to know: \ru{кроссворд}}
\textbf{\ru{кроссворд}} (\emph{krossvórd}) — \textquotedblleft{}a crossword\textquotedblright{}.

\begin{quote}
\ru{Дедушка} \ru{каждое} \ru{утро} \ru{решает} \ru{кроссворд}.
\end{quote}

\begin{quote}
\emph{Dédushka kázhdoye útro resháyet krossvórd.}
\end{quote}

\begin{quote}
Grandad does a crossword every morning.
\end{quote}

Masculine. \emph{\ru{Решать} \ru{кроссворд}} — to do a crossword.

\begin{cousinweb}
\emph{\ru{кроссворд}} was borrowed from English \emph{crossword} (EtymDB).
\end{cousinweb}

\begin{grammarlens}[title={What you've built}]
One more word for telling, explaining and describing at B1.
\end{grammarlens}

\subsection*{Your turn}
\begin{itemize}
\raggedright
  \item \emph{Say it:} \emph{krossvórd}
  \item \emph{Say it:} \emph{krossvórd}, once more
  \item \emph{Say it:} a sentence of your own with \emph{krossvórd}
  \item \emph{Recall:} say the Russian for a collector, then the Russian for a badge, then say \emph{krossvórd} again
\end{itemize}

\begin{tcolorbox}[breakable,colback=teal!4,colframe=teal!35!black,title={Before you move on}]
What does \ru{кроссворд} mean? (\textquotedblleft{}A crossword\textquotedblright{}.) Say it once more.
\end{tcolorbox}
\section[\texorpdfstring{\ru{разгадывать} (razgádyvat)}{разгадывать (razgádyvat)}]{\ru{разгадывать} (razgádyvat) — to solve (a riddle, a puzzle)}

\label{lesson:RU-C503-razgadyvat}



\begin{quote}
Before the new one: say the Russian for a badge, then the Russian for a crossword.
\end{quote}

\subsection*{You'll want to know: \ru{разгадывать}}
\textbf{\ru{разгадывать}} (\emph{razgádyvat}) — \textquotedblleft{}to solve (a riddle, a puzzle)\textquotedblright{}.

\begin{quote}
\ru{Мы} \ru{вместе} \ru{разгадываем} \ru{загадки}.
\end{quote}

\begin{quote}
\emph{My vméste razgádyvayem zagádki.}
\end{quote}

\begin{quote}
We solve riddles together.
\end{quote}

Imperfective; perfective \ru{разгадать}.

\begin{cousinweb}
\emph{\ru{разгадывать}} is built on \emph{\ru{раз}-} and \emph{\ru{гадать}} 'to guess'; EtymDB has no entry for the verb, and \emph{\ru{гадать}} is inherited from Proto-Slavic \emph{gadati} (EtymDB).
\end{cousinweb}

\begin{grammarlens}[title={What you've built}]
One more word for telling, explaining and describing at B1.
\end{grammarlens}

\subsection*{Your turn}
\begin{itemize}
\raggedright
  \item \emph{Say it:} \emph{razgádyvat}
  \item \emph{Say it:} \emph{razgádyvat}, once more
  \item \emph{Say it:} a sentence of your own with \emph{razgádyvat}
  \item \emph{Recall:} say the Russian for a badge, then the Russian for a crossword, then say \emph{razgádyvat} again
\end{itemize}

\begin{tcolorbox}[breakable,colback=teal!4,colframe=teal!35!black,title={Before you move on}]
What does \ru{разгадывать} mean? (\textquotedblleft{}To solve (a riddle, a puzzle)\textquotedblright{}.) Say it once more.
\end{tcolorbox}
